\documentclass[11pt]{article}
% cs250preamble.tex --- shared preamble for CS 250 course notes
%
% This file is \input by main.tex and by each individual module
% file so that each module can still compile standalone. Any
% package or custom-environment change should be made here, not
% in the individual module files.

\usepackage[T1]{fontenc}
\usepackage[margin=1in]{geometry}
\usepackage{amsmath, amssymb, amsthm}
\usepackage{enumitem}
\usepackage{hyperref}
\usepackage{booktabs}
\usepackage{listings}
\usepackage{xcolor}
\usepackage{tikz}
\usetikzlibrary{shapes.geometric, shapes.gates.logic.US, shapes.gates.logic.IEC, arrows.meta, positioning, calc, trees}
\usepackage{bussproofs}
\usepackage{mdframed}

\lstset{
  basicstyle=\ttfamily\small,
  frame=single,
  breaklines=true,
  columns=fullflexible,
  mathescape=true
}

\theoremstyle{definition}
\newtheorem{theorem}{Theorem}[section]
\newtheorem{definition}[theorem]{Definition}
\newtheorem{example}[theorem]{Example}

\hypersetup{
  colorlinks=true,
  linkcolor=blue!60!black,
  urlcolor=blue!60!black
}

% Key-takeaway box: used at the end of major sections and at the end of
% each module to summarise the main points. Expects \item bullets inside.
\newmdenv[
  linewidth=0.8pt,
  linecolor=blue!40!black,
  backgroundcolor=blue!3,
  skipabove=8pt,
  skipbelow=8pt,
  innertopmargin=7pt,
  innerbottommargin=7pt,
  innerleftmargin=9pt,
  innerrightmargin=9pt
]{keytakeawaybox}
\newenvironment{keytakeaway}{%
  \begin{keytakeawaybox}%
  \noindent\textbf{Key takeaways.}%
  \begin{itemize}[leftmargin=1.3em, topsep=2pt, itemsep=1pt, label=\textbullet]%
}{%
  \end{itemize}%
  \end{keytakeawaybox}%
}

% Summary/looking-ahead environment: used once at the end of each module
% to wrap up what the module built and point forward.
\newmdenv[
  linewidth=0.8pt,
  linecolor=gray!60,
  backgroundcolor=gray!5,
  skipabove=8pt,
  skipbelow=8pt,
  innertopmargin=7pt,
  innerbottommargin=7pt,
  innerleftmargin=9pt,
  innerrightmargin=9pt
]{summarybox}

% Sidebar environment: used for tangential asides---historical notes,
% pointers to other areas of math/CS, cross-paradigm remarks---that
% enrich the main narrative without being essential to it. Takes one
% argument: the sidebar's title.
\newmdenv[
  linewidth=0.6pt,
  linecolor=gray!60,
  backgroundcolor=gray!4,
  skipabove=8pt,
  skipbelow=8pt,
  innertopmargin=6pt,
  innerbottommargin=6pt,
  innerleftmargin=9pt,
  innerrightmargin=9pt
]{sidebarbox}
\newenvironment{sidebar}[1]{%
  \begin{sidebarbox}%
  \noindent\textbf{Sidebar: #1.}\par\medskip%
}{%
  \end{sidebarbox}%
}

% Reading-map environment: used at the top of each numbered module to
% indicate Epp coverage, prerequisites, relevant intermezzi, and
% suggested practice problems. Expects four \rmentry{label}{body}
% items inside (or arbitrary paragraphs).
\newmdenv[
  linewidth=0.8pt,
  linecolor=black!55,
  backgroundcolor=yellow!4,
  skipabove=10pt,
  skipbelow=14pt,
  innertopmargin=9pt,
  innerbottommargin=9pt,
  innerleftmargin=11pt,
  innerrightmargin=11pt
]{readingmapbox}
\newcommand{\rmentry}[2]{\par\noindent\textbf{#1}\hspace{0.4em}#2\par}
\newenvironment{readingmap}{%
  \begin{readingmapbox}%
  \noindent{\bfseries\large Reading map}%
  \vspace{2pt}%
  \setlength{\parskip}{4pt plus 1pt}%
}{%
  \end{readingmapbox}%
}


\title{CS 250 --- Cumulative Review 1 (after Module 3)}
\author{}
\date{}

\begin{document}
\maketitle

\section*{The point of this review}

You've now seen three modules' worth of material: sets, functions, and proof-writing basics (Module~1); propositional logic with truth-table semantics (Module~2); and proof systems with natural-deduction-style rules (Module~3). In this short review chapter we pick one small but non-trivial result and prove it \emph{three different ways}, each way using a different module's toolkit. The point is to see explicitly how the three modules connect---how the same mathematical fact can be established semantically (via truth tables), syntactically (via inference rules), or by direct set-theoretic manipulation---and to consolidate the vocabulary before you move on to predicate logic in Module~4.

\section*{The result}

We'll prove the following tautology of propositional logic:
\[
  (p \to q) \;\wedge\; (p \to r) \;\equiv\; p \to (q \wedge r).
\]

In English: ``$p$ implies both $q$ and $r$'' is equivalent to ``$p$ implies the conjunction $q \wedge r$.'' This is a small but useful identity that lets you factor out a common hypothesis across two implications. It also shows up in the proof of every structured induction where you need to prove two things from the same assumption.

\section*{Proof 1: By truth table (Module 2 style)}

Both sides are formulas over the three variables $p, q, r$. Logical equivalence is defined as ``same truth table,'' so we build truth tables for both sides and compare. We need eight rows, one for each assignment to $(p, q, r)$.

\begin{center}
\begin{tabular}{ccc|cc|c|c|c}
\toprule
$p$ & $q$ & $r$ & $p \to q$ & $p \to r$ & $(p \to q) \wedge (p \to r)$ & $q \wedge r$ & $p \to (q \wedge r)$ \\
\midrule
T & T & T & T & T & T & T & T \\
T & T & F & T & F & F & F & F \\
T & F & T & F & T & F & F & F \\
T & F & F & F & F & F & F & F \\
F & T & T & T & T & T & T & T \\
F & T & F & T & T & T & F & T \\
F & F & T & T & T & T & F & T \\
F & F & F & T & T & T & F & T \\
\bottomrule
\end{tabular}
\end{center}

Columns 6 and 8 (the two full sides of the equivalence) match on every row, so the two formulas are logically equivalent. \hfill$\square$

This proof is mechanical and leaves no doubt, but it also gives you no \emph{insight} into why the equivalence holds. If someone asked you to explain the result, you'd say ``it worked out in every row,'' which isn't very satisfying.

\section*{Proof 2: By formal derivation (Module 3 style)}

The same result, established syntactically using the introduction and elimination rules from Module~3. Because this is an equivalence, we need to prove both directions. We'll write the proofs in the numbered-line format from Module~3; you could also do them as Gentzen trees if you prefer.

\subsection*{Direction 1: $(p \to q) \wedge (p \to r) \;\to\; p \to (q \wedge r)$}

\noindent Assume $(p \to q) \wedge (p \to r)$.
\begin{enumerate}
  \item $(p \to q) \wedge (p \to r)$ \hfill (assumption)
  \item $p \to q$ \hfill ($\wedge$-elimination on (1))
  \item $p \to r$ \hfill ($\wedge$-elimination on (1))
  \item Subproof: assuming $p$, show $q \wedge r$:
    \begin{enumerate}[label=(4\alph*)]
      \item $p$ \hfill (assumption)
      \item $q$ \hfill (modus ponens on (4a), (2))
      \item $r$ \hfill (modus ponens on (4a), (3))
      \item $q \wedge r$ \hfill ($\wedge$-introduction on (4b), (4c))
    \end{enumerate}
  \item $p \to (q \wedge r)$ \hfill ($\to$-introduction on subproof (4))
\end{enumerate}

\subsection*{Direction 2: $p \to (q \wedge r) \;\to\; (p \to q) \wedge (p \to r)$}

\noindent Assume $p \to (q \wedge r)$.
\begin{enumerate}
  \item $p \to (q \wedge r)$ \hfill (assumption)
  \item Subproof: assuming $p$, show $q$:
    \begin{enumerate}[label=(2\alph*)]
      \item $p$ \hfill (assumption)
      \item $q \wedge r$ \hfill (modus ponens on (2a), (1))
      \item $q$ \hfill ($\wedge$-elimination on (2b))
    \end{enumerate}
  \item $p \to q$ \hfill ($\to$-introduction on subproof (2))
  \item Subproof: assuming $p$, show $r$:
    \begin{enumerate}[label=(4\alph*)]
      \item $p$ \hfill (assumption)
      \item $q \wedge r$ \hfill (modus ponens on (4a), (1))
      \item $r$ \hfill ($\wedge$-elimination on (4b))
    \end{enumerate}
  \item $p \to r$ \hfill ($\to$-introduction on subproof (4))
  \item $(p \to q) \wedge (p \to r)$ \hfill ($\wedge$-introduction on (3), (5))
\end{enumerate}

Since both directions are derivable, the formulas are equivalent in the proof system. \hfill$\square$

This proof is more verbose than the truth-table version, but it tells you \emph{why} the equivalence holds: you can turn a proof of $(p \to q) \wedge (p \to r)$ into a proof of $p \to (q \wedge r)$ (and back again) by rearranging the same basic steps. In the Curry--Howard view, both formulas correspond to the same ``function of type $p$ returning a pair $(q, r)$,'' which is a genuine computational content that the truth-table proof doesn't touch.

\section*{Proof 3: By sets (Module 1 style)}

Here's a third angle---the one you might not have expected. A predicate on a set is just a subset (the truth set), and propositional equivalences become set identities under this reading. Let's reinterpret the whole claim in set-theoretic language.

Fix any set $U$ and think of $p, q, r$ as subsets of $U$, representing the truth sets of predicates. The connective $p \to q$ is the set of elements where ``if $p$ then $q$'' holds---which is every element where $p$ is false \emph{or} $q$ is true, i.e., $\overline{p} \cup q$. Similarly $q \wedge r$ becomes $q \cap r$. The equivalence we want to prove becomes the set identity
\[
  (\overline{p} \cup q) \cap (\overline{p} \cup r) \;=\; \overline{p} \cup (q \cap r).
\]

This is just \emph{distributivity of union over intersection}, a basic set identity from Module~1! Let's verify it in both directions.

\emph{$(\subseteq)$.} Let $x \in (\overline{p} \cup q) \cap (\overline{p} \cup r)$. Then $x$ is in both $\overline{p} \cup q$ and $\overline{p} \cup r$. Case on whether $x \in \overline{p}$:
\begin{itemize}
  \item If $x \in \overline{p}$, then $x \in \overline{p} \cup (q \cap r)$ by definition of union.
  \item If $x \notin \overline{p}$, then (since $x \in \overline{p} \cup q$) we must have $x \in q$; and (since $x \in \overline{p} \cup r$) we must have $x \in r$. So $x \in q \cap r$, hence $x \in \overline{p} \cup (q \cap r)$.
\end{itemize}
Either way, $x$ is in the right-hand side.

\emph{$(\supseteq)$.} Let $x \in \overline{p} \cup (q \cap r)$. Then either $x \in \overline{p}$ or $x \in q \cap r$.
\begin{itemize}
  \item If $x \in \overline{p}$, then $x \in \overline{p} \cup q$ and $x \in \overline{p} \cup r$, so $x$ is in their intersection.
  \item If $x \in q \cap r$, then $x \in q$ and $x \in r$, so again $x \in \overline{p} \cup q$ and $x \in \overline{p} \cup r$, and $x$ is in their intersection.
\end{itemize}

Both directions hold, so the set identity is established. Translating back to propositional logic: $(p \to q) \wedge (p \to r) \equiv p \to (q \wedge r)$. \hfill$\square$

\section*{What we learned}

Three proofs, three very different flavors:
\begin{itemize}
  \item The \textbf{truth-table proof} is a finite, mechanical check. It's the fastest to write but gives no insight into ``why.'' This is the view from Module~2.
  \item The \textbf{formal derivation} follows the connective structure of the formulas, using one introduction or elimination rule at each step. It's longer but it \emph{explains} the equivalence as a rearrangement of basic proof steps. This is the view from Module~3, and it's the one closest to how a proof assistant like Lean would actually check the fact.
  \item The \textbf{set-theoretic proof} reinterprets the problem entirely, turning propositional equivalence into a set-identity question. It's a reminder that the tools from Module~1 aren't just warm-ups---they're alternative vantage points for everything that follows.
\end{itemize}

The fact that all three proofs exist isn't a coincidence. One of the deepest results about propositional logic---the \emph{completeness theorem}---says that a formula is a semantic tautology (truth-table true) if and only if it has a syntactic derivation (provable from the inference rules). And the set-theoretic reading is what you get when you interpret propositional formulas in the ``Boolean algebra of subsets'' of a fixed universe---a genuine, well-defined algebraic structure.

If you only remember one thing from this review chapter, make it this: \emph{the same fact can wear many different hats}. When you get stuck proving something one way, try translating the problem into a different language. The fact hasn't moved; only your perspective has.

\end{document}
